\documentclass[10pt,a4paper]{amsart}
\usepackage[margin=19mm]{geometry}
\usepackage{amsmath,amssymb,mathtools}
\usepackage[scr=rsfs,cal=euler]{mathalfa}
\usepackage{xfrac}
\usepackage[unicode,hidelinks,bookmarks=false]{hyperref}
\newcommand{\ie}{i.e.\ }
\newcommand{\cf}{cf.\ }
\newcommand{\resp}{respectively\ }
\newcommand{\Subsection}{\subsection}
\providecommand{\nobreakdash}{\nobreak\hbox{-}}
\setlength{\emergencystretch}{2em}
\multlinegap0pt
\usepackage{enumerate}
\usepackage[normalem]{ulem}
\usepackage{amssymb}
\usepackage{tikz}
\newtheorem{thm}{Theorem}
\title{Global regularity and free boundary geometry \\ in the planar Chon\'e-Rochet model}
\author{Shibing Chen, Alessio Figalli,, Yi Ru-Ya Zhang}
\date{Source: arXiv:2603.21196v1 (CC BY 4.0)}
\begin{document}
\maketitle

\noindent\emph{Source and licence.} Global regularity and free boundary geometry \\ in the planar Chon\'e-Rochet model. Version arXiv:2603.21196v1 (CC BY 4.0), distributed by arXiv under the Creative Commons Attribution 4.0 International licence, http://creativecommons.org/licenses/by/4.0/, as recorded in the arXiv licence metadata for this exact version and shown on the abstract page https://arxiv.org/abs/2603.21196v1. Full licence notice: \texttt{licenses/arxiv-2603.21196-CC-BY-4.0.txt}.\par\smallskip

\noindent\textit{Editorial scope.} Counted unit: the article's thm Rigi (source0.tex lines 383--398). The article's complete original proof of it is delivered with it (source0.tex lines 2892--3171, one \texttt{proof} environment of 280 lines, which the article places away from the statement). No mathematics is written, repaired or re-derived: the statement and the proof are the article's own lines, in the article's own notation, and no retained line is substituted or re-flowed.  No further source-local context is needed: every cross-reference of every retained line resolves inside the two delivered spans.  Nothing was omitted from any retained span, so the package carries no omission record.  No retained line contains a citation, so the wrapper carries no bibliography and no bibliography entry.  No retained line contains a figure environment, an image-inclusion command or drawing code, so no image is delivered and the package declares no asset dependency.  Editorial formatting only, in addition: an AMS \texttt{amsart} preamble that restates the article's own declarations and macros used by the retained lines, namely the environments \texttt{thm} on separate counters, together with the standard packages those lines need.  The retained environments print in this document as Theorem 1 in order of appearance, whereas the article prints the same items with the numbering of its own full document.  The complete native source of the article as distributed in the arXiv e-print for this exact version is bundled unchanged as \texttt{sources/196922/source0.tex}.
\begin{thm}\label{Rigi}
Let $u\in C^{1,1}(B_1)$ be a convex function satisfying
\[
u\ge 0 \quad\text{in }B_1\qquad\text{and}
\qquad
u=0 \quad\text{on } S:=\{(x_1,0): -1<x_1<0\}.
\]
Let $a\in(0,\pi/2)$, 
$D:=\{(r,\theta): -\pi+a<\theta<\pi-a\},$
and assume that $\Delta u=1$ in $D\cap B_1$. Then
\[
u(x)=\frac{x_2^2}{2}
\qquad\text{in }D\cap B_1.
\]
In particular, $u$ cannot be strictly convex in $D\cap B_1$.
\end{thm}

\begin{proof}[Proof of Theorem~\ref{Rigi}]
Write
$x=(x_1,x_2)=(r\cos\theta,r\sin\theta).$
Since $u\ge0$ in $B_1$ and $u=0$ on
$S=\{(x_1,0): -1<x_1<0\},$
every point of $S$ is a local minimum. Hence
\begin{equation}\label{eq:grad-on-S}
D u(x_1,0)=0
\qquad\text{for all }-1<x_1\le0.
\end{equation}
Because $u\in C^{1,1}(B_1)$, the gradient is globally Lipschitz: there exists $L\ge1$ such that
\begin{equation}\label{eq:Lip}
|D u(x)-D u(y)|\le L|x-y|
\qquad\forall x,y\in B_1.
\end{equation}
Convexity implies $D^2u\ge0,$
while the identity $\Delta u=1$ in $D$ implies that each first derivative $\partial_i u$
is harmonic in $D$.

\medskip
\noindent\textit{Step 1: a one-sided lower bound for $\partial_1u$.}
Let
\[
y:=\Bigl(-\frac12,0\Bigr)\in S.
\]
Since $u$ is convex, the gradient is monotone:
\[
(D u(x)-D u(y))\cdot (x-y)\ge0.
\]
Using \eqref{eq:grad-on-S}, this gives
\[
\partial_1u(x)\Bigl(x_1+\frac12\Bigr)+\partial_2u(x)\,x_2\ge0.
\]
Thus, whenever $|x_1|\le1/4$,
\[
\partial_1u(x)\ge -4x_2\,\partial_2u(x).
\]
Now let $x\in \partial D\cap B_{1/4}$. Since the sides of $D$ lie in $\{x_1<0\}$, the point
$(x_1,0)$ belongs to $S$, and \eqref{eq:Lip} together with \eqref{eq:grad-on-S} gives
\[
|\partial_2u(x)|
\le |D u(x)-D u(x_1,0)|
\le L|x_2|.
\]
Hence
\[
\partial_1u(x)\ge -4Lx_2^2
\qquad\text{for all }x\in \partial D\cap B_{1/4}.
\]
Set
$C_\ast:=4L.$
Since $u_{11}\ge0$, for each fixed $x_2$ the function
\[
x_1\longmapsto \partial_1u(x_1,x_2)+C_\ast x_2^2
\]
is nondecreasing. Therefore, choosing $r_0=r_0(a)\in(0,1/4]$ so that every horizontal segment
through a point of $D\cap B_{r_0}$ meets $\partial D\cap B_{1/4}$, the previous boundary estimate
propagates into the sector:
\begin{equation}\label{eq:basic-lower}
\partial_1u(x)\ge -C_\ast x_2^2
\qquad\text{for all }x\in D\cap B_{r_0}.
\end{equation}

\medskip
\noindent\textit{Step 2: a harmonic lower barrier for $u_{11}$.}
Fix $a'\in(a,\pi/2)$ and set
\[
D':=\{(r,\theta): -\pi+a'<\theta<\pi-a'\}\subset D,
\qquad
b:=\frac{\pi}{2(\pi-a')}\in\Bigl(\frac12,1\Bigr),
\]
\[
\hat w(r,\theta):=r^b\cos(b\theta).
\]
Then $\hat w$ is harmonic in $D'$, vanishes on $\partial D'$, and is strictly positive in $D'$.
Moreover,
\begin{equation}\label{eq:d1w}
\partial_1\hat w
=
\cos\theta\,\partial_r\hat w-\frac{\sin\theta}{r}\,\partial_\theta\hat w
=
b\,r^{b-1}\cos\bigl((b-1)\theta\bigr),
\qquad
|\partial_1\hat w|\le b\,r^{b-1}.
\end{equation}
Let
$\hat v:=u_{11}.$
Then $\hat v$ is harmonic and nonnegative in $D$.
If $\hat v\equiv0$ in $D$, then $\partial_1u$ is harmonic and independent of $x_1$, hence
\[
\partial_1u(x)=\alpha x_2+\beta
\]
in $D$ for some constants $\alpha,\beta$.
Since $\partial_1u=0$ on $S$, we have $\beta=0$.
The convexity matrix
\[
\begin{pmatrix}
0 & \alpha\\
\alpha & u_{22}
\end{pmatrix}
\]
must be nonnegative semidefinite, so necessarily $\alpha=0$.
Hence $\partial_1u\equiv0$ in $D$.
Since now $\Delta u=u_{22}=1$ in $D$, and $u(0)=|D u(0)|=0$, we conclude that
\[
u(x)=\frac{x_2^2}{2}
\qquad\text{in }D.
\]
So it only remains to exclude the case $\hat v\not\equiv0$.

Assume therefore that $\hat v\not\equiv0$.
By the strong maximum principle, $\hat v>0$ in $D$.
In particular,
\[
m:=\min_{\theta\in[-\pi+a',\,\pi-a']}\hat v(1/2,\theta)>0.
\]
Set
\[
c_0:=m\,2^b,
\qquad
h:=\hat v-c_0\hat w
\quad\text{in }D'\cap B_{1/2}.
\]
Then $h$ is harmonic. On the two sides of $D'$ one has $\hat w=0$, so $h\ge0$ there, while
on the circular arc $\partial B_{1/2}\cap \overline{D'}$,
\[
c_0\hat w\le m\,2^b \Bigl(\frac12\Bigr)^b=m\le \hat v.
\]
Hence $h\ge0$ on $\partial(D'\cap B_{1/2})$, and therefore
\begin{equation}\label{eq:v-lower}
u_{11}(x)=\hat v(x)\ge c_0\hat w(x)
\qquad\text{for all }x\in D'\cap B_{1/2}.
\end{equation}

\medskip
\noindent\textit{Step 3: a nonlinear subharmonic barrier.}
Choose $\gamma\in(1,1/b)$ and define
\[
\delta:=1-b\gamma\in(0,1).
\]
For $\varepsilon>0$, set
\[
\psi_\varepsilon(x):=\varepsilon\,\hat w(x)^\gamma-(C_\ast+1)x_2^2.
\]
Since $\hat w$ is harmonic and positive in $D'$,
\[
\Delta(\hat w^\gamma)
=
\gamma(\gamma-1)\hat w^{\gamma-2}|D \hat w|^2
=
\gamma(\gamma-1)b^2\,r^{b\gamma-2}\cos^{\gamma-2}(b\theta)
\ge \gamma(\gamma-1)b^2\,r^{-(1+\delta)},
\]
because $b\gamma-2=-(1+\delta)$ and $\cos(b\theta)\in(0,1]$ in $D'$.
Since $\Delta(x_2^2)=2$, it follows that
\[
\Delta\psi_\varepsilon
\ge
\varepsilon\,\gamma(\gamma-1)b^2\,r^{-(1+\delta)}-2(C_\ast+1).
\]
Now choose
\begin{equation}\label{eq:r0eps}
\varepsilon:=A\,r_0^{1+\delta},
\qquad
A:=\frac{2(C_\ast+1)}{\gamma(\gamma-1)b^2}.
\end{equation}
Then $\psi_\varepsilon$ is subharmonic in $D'\cap B_{r_0}$.

Set
\[
U:=D'\cap\bigl(\{|x_1|<C_0r_0\}\times \{|x_2|<r_0\}\bigr),
\qquad
C_0:=\cot a'+1,
\]
and define
\[
H:=\partial_1u-\psi_\varepsilon.
\]

\medskip
\noindent\textit{Step 4: $H\ge0$ on $\partial U$.}

\smallskip
\noindent\emph{(a) On the sides of $D'$.}
There $\hat w=0$, so by \eqref{eq:basic-lower},
\[
H=\partial_1u+(C_\ast+1)x_2^2
\ge -C_\ast x_2^2+(C_\ast+1)x_2^2
\ge0.
\]

\smallskip
\noindent\emph{(b) On the top and bottom edges.}
We treat the top edge $\{x_2=r_0\}$, the bottom one being identical.
Set
\[
t_\ell:=-r_0\cot a',
\qquad
t_r:=(\cot a'+1)r_0,
\]
and define
\[
F(t):=H(t,r_0)
\qquad\text{for }t\in[t_\ell,t_r].
\]
Using \eqref{eq:v-lower},
\[
F'(t)
=
u_{11}(t,r_0)-\partial_1\psi_\varepsilon(t,r_0)
\ge c_0\hat w(t,r_0)-\varepsilon\,\partial_t\bigl(\hat w(t,r_0)^\gamma\bigr).
\]
There exist constants $0<c_1\le c_2$, depending only on $a'$, such that
for $t=t_\ell+s$ one has
\begin{equation}\label{eq:lin-growth}
c_1\,r_0^{b-1}s\le \hat w(t,r_0)\le c_2\,r_0^{b-1}s,
\qquad 0\le s\le t_r-t_\ell.
\end{equation}
Also, since $\hat w(t_\ell,r_0)=0$, \eqref{eq:basic-lower} gives
\[
F(t_\ell)=\partial_1u(t_\ell,r_0)+(C_\ast+1)r_0^2\ge r_0^2.
\]
Integrating from $t_\ell$ to $t=t_\ell+s$, and using \eqref{eq:r0eps} and \eqref{eq:lin-growth},
we obtain
\begin{align}
F(t)
&\ge r_0^2 + c_0\int_{t_\ell}^{t}\hat w(\tau,r_0)\,d\tau
-\varepsilon\,\hat w(t,r_0)^\gamma \nonumber\\
&\ge r_0^2+\frac{c_0c_1}{2}\,r_0^{b-1}s^2
-Ac_2^\gamma\,r_0^{1+\delta+\gamma(b-1)}\,s^\gamma \nonumber\\
&= r_0^2+\frac{c_0c_1}{2}\,r_0^{b-1}s^2
-Ac_2^\gamma\,r_0^{2-\gamma}\,s^\gamma. \label{eq:layer-ineq}
\end{align}
Choose $\theta_0>0$ small, depending only on $A,c_2,\gamma$, so that
$F\ge0$ on $[t_\ell,t_\ell+\theta_0r_0]$.
On the complementary interval, the positive term in \eqref{eq:layer-ineq}
is of order $r_0^{1+b}$ while the negative term is of order $r_0^2$;
since $1+b<2$, the positive contribution dominates for $r_0$ sufficiently small.
Hence $F\ge0$ on the whole top edge.

\smallskip
\noindent\emph{(c) On the right side.}
Fix $x_2\in[-r_0,r_0]$ and let
\[
t_\ell(x_2):=-|x_2|\cot a'
\]
be the left endpoint of the horizontal section of $D'$ at height $x_2$.
Starting from the boundary value $H(t_\ell(x_2),x_2)\ge0$ given by part (a), and repeating
the argument from part (b) along the horizontal segment from $t_\ell(x_2)$ to $t_r$, we obtain
\[
H(t_r,x_2)\ge0
\]
for all $|x_2|\le r_0$.
Therefore $H\ge0$ on $\partial U$.

\medskip
\noindent\textit{Step 5: interior comparison and contradiction.}
Inside $U$, the function $\partial_1u$ is harmonic while $\psi_\varepsilon$ is subharmonic,
so $H=\partial_1u-\psi_\varepsilon$ is superharmonic.
Thus, by the minimum principle and Step~4 we deduce that
$H\ge0$ inside $U$, and therefore
\[
\partial_1u\ge \psi_\varepsilon
\qquad\text{in }U.
\]
In particular, for $x=(x_1,0)$ with $x_1>0$ sufficiently small,
\[
\partial_1u(x_1,0)\ge \varepsilon\,\hat w(x_1,0)^\gamma
=
\varepsilon\,x_1^{b\gamma}.
\]
Since $b\gamma<1$, this contradicts the Lipschitz continuity of $\partial_1u$ at the origin
together with $\partial_1u(0)=0$.
Therefore the alternative $\hat v\not\equiv0$ is impossible, and we are in the rigid case
\[
u(x)=\frac{x_2^2}{2}
\qquad\text{in }D\cap B_1.
\]
In particular, $u$ cannot be strictly convex in $D\cap B_1$.
\end{proof}

\end{document}
